% !TeX program = xelatex
% !TeX encoding = UTF-8
\documentclass[12pt,a4paper]{article}

\usepackage[margin=2.5cm]{geometry}
\usepackage{amsmath,amssymb,amsthm}
\usepackage{xeCJK}
\usepackage{graphicx}
\usepackage[hidelinks]{hyperref}
\usepackage{fancyhdr}
\usepackage{needspace}

\setCJKmainfont{Microsoft JhengHei}

\makeatletter
\newcommand{\totalpages}{\@abspage@last}
\makeatother
\pagestyle{fancy}
\fancyhf{}
\renewcommand{\headrulewidth}{0pt}
\fancyfoot[C]{\thepage/\totalpages}
\fancypagestyle{plain}{%
  \fancyhf{}%
  \renewcommand{\headrulewidth}{0pt}%
  \fancyfoot[C]{\thepage/\totalpages}%
}

\setlength{\parindent}{0pt}
\setlength{\parskip}{0.5em}
\raggedbottom
\title{Linear System --- Homework 5}
\author{葉彥辰\quad P46154345}
\date{}

\begin{document}
\maketitle
\begin{center}Fall 2026\quad Due: Friday, Oct 31\end{center}

\section*{Exercise 1}
Suppose $A$ is a $3\times 3$ matrix and
\[
\det(sI-A)=s^3+2s^2+s+1.
\]

\needspace{6cm}
\textbf{(a)} Express $A^3$ in terms of $I$, $A$, and $A^2$.
\subsection*{Solution}
% 在這裡撰寫解答。

\begin{align*}
	p(s) &= \det (sI - A) = s^3 + 2s^2 + s + 1 \\
	p(A) &= A^3 + 2A^2 + A + I  = 0\\
	\Rightarrow A^3 &= -2A^2 -A-I
\end{align*}

\textbf{(b)} Express $A^5$ in terms of $I$, $A$, and $A^2$.
\subsection*{Solution}
% 在這裡撰寫解答。

\[
A^5 = A^2 A^3 = -2A^4 - A^3 - A^2
\]
\begin{align*}
A^4 &= A^3A = -2A^3 - A^2 - A  \\
&= -2( -2A^2 -A-I) - A^2 - A \\
&= 4A^2 + 2A + 2I -A^2 - A \\
&=3A^2 + A + 2I
\end{align*}

\begin{align*}
A^5 &= -2 (3A^2 + A +2I) - (-2A^2 - A - I) - A^2 \\
    &= -6A^2 -2A -4I + 2A^2 +A + I - A^2 \\
    &= -5A^2-A-3I.
\end{align*}
\vspace{0.5cm}


\textbf{(c)} Express $A^{-1}$ in terms of $I$, $A$, and $A^2$.
\subsection*{Solution}
% 在這裡撰寫解答。
Since $p(0)=\det(-A)=1$ and $A$ is $3\times3$, $\det A=-1\ne0$.
Thus $A$ is invertible.
\begin{align*}
	A^{-1}p(A) &= A^2 + 2A + I + A^{-1} = 0 \\
	\Rightarrow A^{-1} &= -A^2 - 2A - I
\end{align*}

\section*{Exercise 2}
Without doing any matrix multiplications, compute $A^4$ for
\[
A=\begin{bmatrix}
0&1&0&0\\
0&0&1&0\\
0&0&0&1\\
564&0&0&0
\end{bmatrix}.
\]
Justify your answer.
\subsection*{Solution}
% 在這裡撰寫解答。
\begin{align*}
p(s)&=\det(sI-A)=
\begin{vmatrix}
s&-1&0&0\\
0&s&-1&0\\
0&0&s&-1\\
-564&0&0&s
\end{vmatrix} \\
&= s \begin{vmatrix}
	s & -1 & 0 \\
	0 & s & -1 \\
	0 & 0 & s \\
\end{vmatrix} - (-564)\begin{vmatrix}
-1 &  0 & 0 \\
s & -1 & 0 \\
0 & s & -1
\end{vmatrix} \\
&= s^4 - 564 \\ \\
\Rightarrow p(A) &= A^4 - 564I = 0 \quad \text{(Cayley--Hamilton theorem)} \\
\therefore A^4 &= 564I
\end{align*}


\section*{Exercise 3}
Consider the differential equation
\[
\dot{x}=Ax,
\]
where $A$ is a square matrix. Show that if $A$ has $j2\pi$ as an eigenvalue, then there is a nonzero initial state $x_0$ such that the solution $x$ satisfies
\[
x(1)=x(0)=x_0.
\]
\subsection*{Solution}
% 在這裡撰寫解答。
Since $j2\pi$ is an eigenvalue of $A$, choose a nonzero eigenvector
$v$ such that $Av=j2\pi v$. Set $x_0=v$ and let
\[
x(t)=e^{j2\pi t}v.
\]
Then $x(0)=v=x_0$ and
\[
\dot{x}(t)=j2\pi e^{j2\pi t}v
=e^{j2\pi t}Av=Ax(t).
\]
\[
\Rightarrow x(1) = e^{j2\pi}x_0 = x_0 = x(0)
\]
If $A$ is real and a real initial state is required, take the real
or imaginary part of this solution whose initial value is nonzero.
Both parts satisfy $\dot{x}=Ax$ and return to their initial values at $t=1$.


\section*{Exercise 4}
Compute $e^{At}$ at $t=\ln(2)$ for
\[
A=\begin{bmatrix}0&1\\-1&0\end{bmatrix}
\]
by:

\needspace{6cm}
\textbf{(a)} use Jordan form using eigenvalues and eigenvectors of $A$;
\subsection*{Solution}
% 在這裡撰寫解答。
\begin{align*}
	\det (sI -A) = \begin{vmatrix}
		s & -1 \\ 1 & s
	\end{vmatrix} = s^2 + 1 = 0, \quad \therefore s = \pm j
\end{align*}
for $s = j$

\[
\begin{bmatrix}
	j & -1 \\ 1 & j
\end{bmatrix} v = \begin{bmatrix}
	j & -1 \\ 1 & j
\end{bmatrix}\begin{bmatrix}
v_1 \\ v_2
\end{bmatrix} = \begin{bmatrix}
0 \\ 0
\end{bmatrix} \quad \Rightarrow \begin{bmatrix}
v_1 \\ v_2 
\end{bmatrix} = \begin{bmatrix}
j \\ -1
\end{bmatrix}
\]

for $s = -j$
\[
\begin{bmatrix}
	-j & -1 \\ 1& -j 
\end{bmatrix} v = \begin{bmatrix}
-j & -1 \\ 1& -j 
\end{bmatrix}\begin{bmatrix}
v_1 \\ v_2
\end{bmatrix} = \begin{bmatrix}
0 \\ 0
\end{bmatrix} \quad \begin{bmatrix}
v_1 \\ v_2
\end{bmatrix} = \begin{bmatrix}
j \\ 1
\end{bmatrix}
\]

Therefore 
\[
T = \begin{bmatrix}
	j & j \\ -1 & 1
\end{bmatrix}, \quad J = \begin{bmatrix}
j & 0 \\ 0 & -j
\end{bmatrix}, \quad T^{-1} = {1\over 2j} \begin{bmatrix}
1 & -j \\ 1 & j
\end{bmatrix}
\]

By $A = TJT^{-1}$
\[
e^{At} = T \begin{bmatrix}
	e^{jt} & 0 \\ 0 & e^{-jt} 
\end{bmatrix} T^{-1}
\]
Hence, at $ t = \ln 2$ 
\begin{align*}
e^{A\ln2} &= T\begin{bmatrix}
	e^{j\ln2} & 0 \\ 0 & e^{-j\ln2}
\end{bmatrix}T^{-1} = {1\over 2j}\begin{bmatrix}
		j & j \\ -1 & 1
\end{bmatrix}\begin{bmatrix}
2^j & 0 \\ 0 & 2^{-j}
\end{bmatrix} \begin{bmatrix}
1 & -j \\ 1 & j
\end{bmatrix} \\&= {1\over2}\begin{bmatrix}
2^{-j} + 2^j & -j(2^j - 2^{-j}) \\
j(2^j - 2^{-j}) &2^{-j} + 2^j
\end{bmatrix} = \begin{bmatrix}
\cos(\ln2) & \sin(\ln2) \\ -\sin(\ln2) & \cos(\ln2)
\end{bmatrix}
\end{align*}

\vspace{3cm}

\needspace{6cm}
\textbf{(b)} using the Laplace transform.
\subsection*{Solution}
% 在這裡撰寫解答。

By
\[
\mathcal L \{ e^{At}\} = (sI - A)^{-1} \quad \Leftrightarrow \quad e^{At} = \mathcal L^{-1} \{(sI - A)^{-1}\}
\]
\begin{align*}
	(sI - A)^{-1} &= \begin{bmatrix}
		s & -1 \\ 1 & s
	\end{bmatrix}^{-1} \\
	&=  \begin{bmatrix}
		{s\over s^2 + 1} & {1\over s^2 + 1} \\ -{1\over s^2 + 1} & {s\over s^2 + 1}
	\end{bmatrix} \\
	\Rightarrow \mathcal L^{-1} \{ (sI -A)^{-1}\} = e^{At} &= \mathcal L^{-1} \left\{ \begin{bmatrix}
		{s\over s^2 + 1} & {1\over s^2 + 1} \\ -{1\over s^2 + 1} & {s\over s^2 + 1}
	\end{bmatrix} \right\} = \begin{bmatrix}
	\cos t & \sin t \\ -\sin t & \cos t
	\end{bmatrix}
\end{align*}

Hence, at $t = \ln2$	
\[
e^{At} = \begin{bmatrix}
	\cos(\ln2) & \sin(\ln2) \\ -\sin(\ln2) & \cos(\ln2)
\end{bmatrix}
\]


\section*{Exercise 5}
Compute $e^{At}$ for
\[
A=\begin{bmatrix}3&-5\\-5&3\end{bmatrix}
\]
by:

\textbf{(a)} use Jordan form using eigenvalues and eigenvectors of $A$;
\subsection*{Solution}
% 在這裡撰寫解答。

\[
\det (sI - A) = \det \begin{bmatrix}
	s -3 & 5 \\ 5 & s - 3
\end{bmatrix} = s^2-6s-16 = (s-8)(s+2)
\]
therefore, eigenvalues $s=8$ and $s=-2$

for $s = 8$
\begin{align*}
	(sI - A)v &= \begin{bmatrix}
		5 & 5 \\ 5 & 5
	\end{bmatrix}v = \begin{bmatrix}
	0 \\ 0
	\end{bmatrix}, \quad v = \begin{bmatrix}
	1 \\ -1 
	\end{bmatrix}
\end{align*}

for $s = -2$
\begin{align*}
	(sI - A)v &= \begin{bmatrix}
		-5 & 5 \\ 5 & -5
	\end{bmatrix}	v = \begin{bmatrix}
	0 \\ 0
	\end{bmatrix}, \quad v = \begin{bmatrix}
	1 \\ 1 
\end{bmatrix}
\end{align*}

\[
T = \begin{bmatrix}
	1 & 1 \\ -1 & 1	
\end{bmatrix},\quad J = \begin{bmatrix}
8 & 0 \\ 0 & -2
\end{bmatrix}, \quad T^{-1} = {1\over2}\begin{bmatrix}
1 & -1 \\ 1 & 1
\end{bmatrix}
\]
Finally
\begin{align*}
	e^{At} &= {1\over2} \begin{bmatrix}
		1 & 1 \\ -1 & 1	
	\end{bmatrix}\begin{bmatrix}
	e^{8t} & 0 \\ 0 & e^{-2t}
	\end{bmatrix}\begin{bmatrix}
	1 & -1 \\ 1 & 1
	\end{bmatrix} \\
	&= {1\over2}\begin{bmatrix}
		e^{8t} & e^{-2t} \\
		-e^{8t} & e^{-2t}
	\end{bmatrix}\begin{bmatrix}
	1 & -1 \\ 1 & 1
	\end{bmatrix}\\
	&= {1\over2 } \begin{bmatrix}
		e^{8t} + e^{-2t} & -e^{8t} + e^{-2t} \\
		-e^{8t} + e^{-2t} & e^{8t} + e^{-2t}
	\end{bmatrix} \\
    &= {e^{3t}\over 2} \begin{bmatrix}
        e^{5t} + e^{-5t} & -(e^{5t} - e^{-5t}) \\
        -(e^{5t} - e^{-5t}) & e^{5t} + e^{-5t}
    \end{bmatrix} \\
    &= e^{3t} \begin{bmatrix}
        \cosh(5t) & -\sinh(5t) \\
        -\sinh(5t) & \cosh(5t)
    \end{bmatrix}
\end{align*}

\textbf{(b)} using the Laplace transform.
\subsection*{Solution}
% 在這裡撰寫解答。
By
\[
\mathcal L \{e^{At}\} = (sI -A)^{-1} \Rightarrow e^{At} = \mathcal{L}^{-1} \{ (sI - A)^{-1} \}
\]

\begin{align*}
    (sI-A)^{-1}
    &= \frac{1}{(s-8)(s+2)}
    \begin{bmatrix}
        s-3 & -5 \\ -5 & s-3
    \end{bmatrix}.
\end{align*}

For the $(1,1)$ entry:
\[
\frac{s-3}{(s-8)(s+2)}
=\frac{1}{2}\frac{1}{s-8}+\frac{1}{2}\frac{1}{s+2}.
\]

For the $(1,2)$ entry:
\[
\frac{-5}{(s-8)(s+2)}
=-\frac{1}{2}\frac{1}{s-8}+\frac{1}{2}\frac{1}{s+2}.
\]

For the $(2,1)$ entry:
\[
\frac{-5}{(s-8)(s+2)}
=-\frac{1}{2}\frac{1}{s-8}+\frac{1}{2}\frac{1}{s+2}.
\]

For the $(2,2)$ entry:
\[
\frac{s-3}{(s-8)(s+2)}
=\frac{1}{2}\frac{1}{s-8}+\frac{1}{2}\frac{1}{s+2}.
\]

\begin{align*}
\Rightarrow \mathcal{L}^{-1} \{(sI - A)^{-1}\} &= {1\over2} \begin{bmatrix}
		e^{8t} + e^{-2t} & -e^{8t} + e^{-2t} \\
	-e^{8t} + e^{-2t} & e^{8t} + e^{-2t}
\end{bmatrix}\\&= e^{3t} \begin{bmatrix}
\cosh(5t) & -\sinh(5t) \\
-\sinh(5t) & \cosh(5t)
\end{bmatrix}
\end{align*}
\vspace{3cm}


\needspace{11cm}
\section*{Exercise 6}
Compute $e^{At}$ at $t=\ln(2)$ for
\[
A=\begin{bmatrix}0&1\\-1&0\end{bmatrix}.
\]
\subsection*{Solution}
% 在這裡撰寫解答。
\begin{align*}
	\det (sI -A) = \begin{vmatrix}
		s & -1 \\ 1 & s
	\end{vmatrix} = s^2 + 1 = 0, \quad \therefore s = \pm j
\end{align*}
for $s = j$

\[
\begin{bmatrix}
	j & -1 \\ 1 & j
\end{bmatrix} v = \begin{bmatrix}
	j & -1 \\ 1 & j
\end{bmatrix}\begin{bmatrix}
	v_1 \\ v_2
\end{bmatrix} = \begin{bmatrix}
	0 \\ 0
\end{bmatrix} \quad \Rightarrow \begin{bmatrix}
	v_1 \\ v_2 
\end{bmatrix} = \begin{bmatrix}
	j \\ -1
\end{bmatrix}
\]

for $s = -j$
\[
\begin{bmatrix}
	-j & -1 \\ 1& -j 
\end{bmatrix} v = \begin{bmatrix}
	-j & -1 \\ 1& -j 
\end{bmatrix}\begin{bmatrix}
	v_1 \\ v_2
\end{bmatrix} = \begin{bmatrix}
	0 \\ 0
\end{bmatrix} \quad \begin{bmatrix}
	v_1 \\ v_2
\end{bmatrix} = \begin{bmatrix}
	j \\ 1
\end{bmatrix}
\]

Therefore 
\[
T = \begin{bmatrix}
	j & j \\ -1 & 1
\end{bmatrix}, \quad J = \begin{bmatrix}
	j & 0 \\ 0 & -j
\end{bmatrix}, \quad T^{-1} = {1\over 2j} \begin{bmatrix}
	1 & -j \\ 1 & j
\end{bmatrix}
\]

By $A = TJT^{-1}$
\[
e^{At} = T \begin{bmatrix}
	e^{jt} & 0 \\ 0 & e^{-jt} 
\end{bmatrix} T^{-1}
\]
Hence, at $ t = \ln 2$ 
\begin{align*}
	e^{A\ln2} &= T\begin{bmatrix}
		e^{j\ln2} & 0 \\ 0 & e^{-j\ln2}
	\end{bmatrix}T^{-1} = {1\over 2j}\begin{bmatrix}
		j & j \\ -1 & 1
	\end{bmatrix}\begin{bmatrix}
		2^j & 0 \\ 0 & 2^{-j}
	\end{bmatrix} \begin{bmatrix}
		1 & -j \\ 1 & j
	\end{bmatrix} \\&= {1\over2}\begin{bmatrix}
		2^{-j} + 2^j & -j(2^j - 2^{-j}) \\
		j(2^j - 2^{-j}) &2^{-j} + 2^j
	\end{bmatrix} = \begin{bmatrix}
		\cos(\ln2) & \sin(\ln2) \\ -\sin(\ln2) & \cos(\ln2)
	\end{bmatrix}
\end{align*}

\vspace{3cm}

\needspace{11cm}
\section*{Exercise 7}
Obtain (by hand) the state response and output response $x(\cdot)$ of the following system due to a unit impulse input and the zero initial conditions. Determine whether the state response contains all the system modes.
\begin{align*}
\dot{x}_1&=-5x_1+2x_2+u,\\
\dot{x}_2&=-12x_1+5x_2+u,\\
y&=2x_1+3x_2.
\end{align*}
\subsection*{Solution}
% 在這裡撰寫解答。
\[
\begin{cases}
	\begin{bmatrix}
		\dot x_1 \\ \dot x_2
	\end{bmatrix} &= \begin{bmatrix}
		-5 & 2 \\ -12 & 5
	\end{bmatrix}\begin{bmatrix}
	x_1 \\ x_2
	\end{bmatrix} + \begin{bmatrix}
	1 \\ 1
	\end{bmatrix}u = Ax + Bu \\ \\
	y &= \begin{bmatrix}
		2 & 3 
	\end{bmatrix}
	\begin{bmatrix}
		x_1 \\ x_2
	\end{bmatrix} = Cx + Du
\end{cases} 
\]
\[
\dot x - Ax(t) = Bu(t)
\]
Multiply both sides on the left by $e^{-At}$. Since $A$ commutes with $e^{-At}$,
\[
e^{-At}\dot x(t) - e^{-At}Ax(t) = e^{-At}Bu(t) \Rightarrow{d\over dt} (e^{-At}x(t)) = e^{-At}Bu(t)
\]
Integrate from $0^-$ to $t$:
\[
e^{-At}x(t) - x(0^-) = \int_{0^-}^t e^{-A\tau}Bu(\tau)d\tau
\]
Therefore
\[
x(t) = e^{At}x(0^-)+ e^{At}\int_{0^-}^t e^{-A\tau}Bu(\tau)d\tau 
\]
For zero initial conditions, $x(0^-)=0$, and a unit impulse input, $u(t)=\delta(t)$,

\begin{align*}
	x(t) &= e^{At}\int_{0^-}^t e^{-A\tau}B\delta(\tau)d\tau \\
	&=e^{At} B, \quad t > 0
\end{align*}
and
\[
y = Cx \Rightarrow y = Ce^{At}B, \quad t > 0.
\]

\vspace{1cm}
Solving $e^{At}$

\[
A = \begin{bmatrix}
	-5 & 2 \\ -12 & 5
\end{bmatrix}
\]

\begin{align*}
	\det (sI - A) &= \begin{vmatrix}
		s + 5 & -2 \\ 12 & s - 5 
	\end{vmatrix} \\
	&= s^2 - 25 + 24 = s^2 -1  = (s + 1)(s - 1) = 0
\end{align*}

Hence, $A$ has two eigenvalues, $1$ and $-1$.

for $s = 1$

\begin{align*}
	(sI - A)v = \begin{bmatrix}
		6 & -2 \\ 12 & -4
	\end{bmatrix}v = 0 \Rightarrow v = \begin{bmatrix}
	1 \\ 3
	\end{bmatrix}
\end{align*}

for $ s= -1$

\[
(sI - A)v = \begin{bmatrix}
	4 & -2 \\ 12 & -6
\end{bmatrix}v = 0 \Rightarrow v=\begin{bmatrix}
1 \\ 2
\end{bmatrix}
\]

Therefore
\[
T = \begin{bmatrix}
	1 & 1 \\
	3 & 2
\end{bmatrix}, \quad J = \begin{bmatrix}
1 & 0 \\ 0 & -1
\end{bmatrix}, \quad T^{-1} = {1\over -1} \begin{bmatrix}
2 & -1 \\ -3 & 1
\end{bmatrix} = \begin{bmatrix}
-2 & 1 \\ 3 & -1
\end{bmatrix}
\]

\begin{align*}
e^{At} &= T e^{Jt} T^{-1} = \begin{bmatrix}
	1 & 1 \\ 3 & 2 
\end{bmatrix}	\begin{bmatrix}
e^{t} & 0 \\ 0 & e^{-t}
\end{bmatrix}\begin{bmatrix}
-2 & 1 \\ 3 & -1
\end{bmatrix} \\
&= \begin{bmatrix}
	e^{t} & e^{-t} \\
	3e^{t} & 2e^{-t}
\end{bmatrix}\begin{bmatrix}
-2 & 1 \\ 3 & -1
\end{bmatrix} \\
&= \begin{bmatrix}
	-2e^{t} + 3e^{-t} & e^{t} -e^{-t} \\
	-6e^{t} + 6e^{-t} & 3e^{t}-2e^{-t}
\end{bmatrix}
\end{align*}

Finally, 
\begin{align*}
 x(t) &= e^{At}B \\
 &= \begin{bmatrix}
	-2e^{t} + 3e^{-t} & e^{t} -e^{-t} \\
-6e^{t} + 6e^{-t} & 3e^{t}-2e^{-t}
 \end{bmatrix}	\begin{bmatrix}
 1 \\ 1
 \end{bmatrix} \\
 &= \begin{bmatrix}
 (-2e^t+3e^{-t})+(e^t-e^{-t}) \\
 (-6e^t+6e^{-t})+(3e^t-2e^{-t})
 \end{bmatrix} \\
 &= \begin{bmatrix}
 -e^t+2e^{-t} \\
 -3e^t+4e^{-t}
 \end{bmatrix}, \qquad t>0.
\end{align*}
and 
\begin{align*}
y(t) &= Ce^{At}B \\
&= \begin{bmatrix}2 & 3\end{bmatrix}
\begin{bmatrix}
-e^t+2e^{-t} \\
-3e^t+4e^{-t}
\end{bmatrix} \\
&= 2(-e^t+2e^{-t})+3(-3e^t+4e^{-t}) \\
&= -11e^t+16e^{-t}, \qquad t>0.
\end{align*}

The state response can be written as
\[
x(t)=-\begin{bmatrix}1\\3\end{bmatrix}e^t
     +2\begin{bmatrix}1\\2\end{bmatrix}e^{-t}.
\]
Both modes $e^t$ and $e^{-t}$, corresponding to the eigenvalues $1$ and $-1$,
have nonzero coefficients. Therefore, the state response contains all the system modes.


\vspace{4cm}

\end{document}
